\documentclass[11pt]{article}
\usepackage{amsmath,amssymb,amsthm,geometry}
\geometry{margin=1in}
\usepackage{cite}

\newcommand{\R}{\mathbb{R}}
\newcommand{\E}{\mathbb{E}}
\newcommand{\W}{\mathbb{W}}
\newcommand{\one}{\mathbf{1}}
\newcommand{\dd}{\mathrm{d}}

\newtheorem{theorem}{Theorem}
\newtheorem{lemma}[theorem]{Lemma}
\newtheorem{corollary}[theorem]{Corollary}
\theoremstyle{remark}
\newtheorem{remark}[theorem]{Remark}

\title{Minimax rate for a nonlinear statistical inverse problem for Darcy's equation}
\author{}
\date{}

\begin{document}
\maketitle

\begin{abstract}
We determine the minimax rate for estimating the conductivity \(f\) in the elliptic divergence-form equation
\(\operatorname{div}(f\nabla u)=1\), \(u|_{\partial D}=0\), from a Gaussian white-noise observation of \(u\). 
For \(f\) in the Hölder--Zygmund class \(C^\alpha(D)\), \(\alpha>1\), and sufficiently close to the constant \(1\), the rate in \(L^2(D)\)-mean squared error is
\[
n^{-\frac{2\alpha}{2\alpha+d+2}},
\]
up to constants. The proof combines a local linearization of the conductivity-to-solution map with the spectral asymptotics of the linearized pseudodifferential operator of order \(-1\).
\end{abstract}

\section{Main result}

Let \(D\subset\R^d\), \(d\ge2\), be the Euclidean unit ball. For \(\alpha>1\) and sufficiently small \(r>0\), define
\[
\mathcal F(\alpha,r)=\{h\in C^\alpha(D):\ \|h-\one\|_{C^\alpha(D)}\le r\},
\]
where \(\one\) is the constant function equal to \(1\). For \(f\in\mathcal F(\alpha,r)\), let \(u_f\) be the unique solution of
\[
\operatorname{div}(f\nabla u)=1\quad\text{in }D,\qquad u|_{\partial D}=0.
\]
We observe
\[
Y_n=u_f+\frac1{\sqrt n}\mathbb W,
\]
where \(\mathbb W\) is a Gaussian white noise on \(L^2(D)\).

\begin{theorem}\label{thm:main}
Let \(d\ge2\) and \(\alpha>1\). For every sufficiently small fixed \(r>0\), there exist constants \(c=c(d,\alpha,r)\ge1\) and \(n_0\) such that for all \(n\ge n_0\),
\[
c^{-1}n^{-\frac{2\alpha}{2\alpha+d+2}}
\le
\inf_{\hat f_n}\sup_{f_0\in\mathcal F(\alpha,r)}
\E_{f_0}\bigl[\|\hat f_n-f_0\|_{L^2(D)}^2\bigr]
\le
c\,n^{-\frac{2\alpha}{2\alpha+d+2}}.
\]
In particular, the minimax rate is
\[
n^{-\frac{2\alpha}{2\alpha+d+2}}
\]
up to multiplicative constants, and no logarithmic factor is needed.
\end{theorem}

The rest of the paper is devoted to the proof of Theorem~\ref{thm:main}.

\section{Linearization of the inverse problem}

Let \(g=f-\one\). For \(f\in\mathcal F(\alpha,r)\), we have \(\|g\|_{C^\alpha(D)}\le r\). Let \(u_0\) be the solution for \(f=\one\):
\[
\Delta u_0=1\quad\text{in }D,\qquad u_0|_{\partial D}=0.
\]
Since \(D\) is the unit ball,
\[
u_0(x)=\frac{|x|^2-1}{2d},\qquad \nabla u_0(x)=\frac{x}{d}.
\]
Put
\[
a(x)=\nabla u_0(x)=\frac{x}{d}.
\]
Let \(v=u_f-u_0\). Subtracting the equation for \(u_0\) from the equation for \(u_f\), we get
\begin{equation}\label{eq:nonlinear-v}
\Delta v=-\operatorname{div}(a g)-\operatorname{div}(g\nabla v),\qquad v|_{\partial D}=0.
\end{equation}
Define the linear operator \(A\) by
\begin{equation}\label{eq:A}
\Delta(A g)=-\operatorname{div}(a g),\qquad (A g)|_{\partial D}=0.
\end{equation}
Then \eqref{eq:nonlinear-v} can be written as
\[
v=A g+R(g),
\]
where
\[
\Delta R(g)=-\operatorname{div}(g\nabla v),\qquad R(g)|_{\partial D}=0.
\]
By elliptic regularity and the smallness of \(g\) in \(C^\alpha\), the remainder \(R(g)\) is a higher-order term. More precisely, for \(\|g\|_{C^\alpha(D)}\le r\) with \(r\) small,
\[
\|R(g)\|_{H^{\alpha+1}(D)}
\le C_\alpha r\|g\|_{H^\alpha(D)}.
\]
Consequently, the map \(g\mapsto u_f-u_0\) is a \(C^2\) perturbation of the linear map \(A\) with second derivative bounded on \(\mathcal F(\alpha,r)\). The local linearization lemma for nonlinear statistical inverse problems with Gaussian white noise \cite[Lemma 2.3]{nickl2019} then implies that the minimax risk in the nonlinear model is, up to constants, the same as that in the linear Gaussian model
\begin{equation}\label{eq:linear-model}
Y=A g+\frac1{\sqrt n}\mathbb W,\qquad g\in B_{C^\alpha}(0,r).
\end{equation}
Thus it remains to determine the minimax rate for \eqref{eq:linear-model}.

\section{Spectral asymptotics of the linearized operator}

\begin{lemma}\label{lem:spectral}
The operator \(A\) defined by \eqref{eq:A} is a classical pseudodifferential operator of order \(-1\) on \(D\). Its principal symbol is
\[
\sigma_A(x,\xi)=\frac{i\,x\cdot\xi}{d|\xi|^2},\qquad (x,\xi)\in T^*D.
\]
In particular, its singular values \(s_k(A)\), \(k=1,2,\dots\), satisfy
\[
s_k(A)\asymp k^{-1/d}.
\]
Equivalently, the counting function \(N(t)=\#\{k:\ s_k(A)>t\}\) satisfies
\[
N(t)\asymp t^{-d},\qquad t\to0^+.
\]
\end{lemma}

\begin{proof}
The operator \(\Delta^{-1}\) has symbol \(-|\xi|^{-2}\), and the operator \(\operatorname{div}(a\,\cdot)\) has principal symbol \(i\,a(x)\cdot\xi\). Hence
\[
\sigma_A(x,\xi)=-\bigl(-|\xi|^{-2}\bigr)\bigl(i\,a(x)\cdot\xi\bigr)
=\frac{i\,x\cdot\xi}{d|\xi|^2},
\]
which is homogeneous of degree \(-1\) in \(\xi\). 

The principal symbol of \(A^*A\) is
\[
|\sigma_A(x,\xi)|^2=\frac{|x\cdot\xi|^2}{d^2|\xi|^4}.
\]
This is an order \(-2\) symbol. Although it vanishes on the measure-zero set \(\{x=0\}\cup\{\xi\cdot x=0\}\), a standard Weyl law for pseudodifferential operators with a principal symbol of this type gives
\[
\#\{\lambda_j(A^*A)>\lambda\}\asymp \lambda^{-d/2},\qquad \lambda\to0^+.
\]
Indeed, the phase-space volume of \(\{(x,\xi):|\sigma_A(x,\xi)|^2>\lambda\}\) is asymptotic to \(\lambda^{-d/2}\). Therefore the eigenvalues \(\lambda_j(A^*A)\) satisfy \(\lambda_j\asymp j^{-2/d}\), and hence
\[
s_j(A)=\sqrt{\lambda_j}\asymp j^{-1/d}.
\]
See \cite[Chapter 29]{hormander1985} for the relevant Weyl law.
\end{proof}

\section{Minimax rate for the linear problem}

We now prove the linear minimax rate. For convenience, we use a wavelet-packet description. Since \(A\) is an order \(-1\) pseudodifferential operator, at frequency scale \(M\) it behaves like multiplication by \(M^{-1}\) on wavelet packets concentrated away from the zero set of \(a\). The zero set of \(a\) is only the point \(x=0\), and wavelets near \(x=0\) affect only \(O(1)\) degrees of freedom per scale; hence they do not change the rate. This is the standard localization argument for pseudodifferential inverse problems.

\begin{lemma}\label{lem:linear-rate}
Consider the Gaussian white-noise model
\[
Y=A g+\frac1{\sqrt n}\mathbb W,
\]
where \(g\in H^\alpha(D)\) with \(\|g\|_{H^\alpha(D)}\le1\). Then the minimax \(L^2(D)\)-risk satisfies
\[
\inf_{\hat g}\sup_{\|g\|_{H^\alpha}\le1}
\E\|\hat g-g\|_{L^2(D)}^2
\asymp n^{-\frac{2\alpha}{2\alpha+d+2}}.
\]
\end{lemma}

\begin{proof}
We give both bounds.

\medskip
\noindent\textbf{Upper bound.}
Let \(\{e_k\}\) be an orthonormal basis of \(A^*A\), and let \(A e_k=s_k f_k\) with \(s_k\asymp k^{-1/d}\). Write
\[
g=\sum_{k\ge1}\theta_k e_k.
\]
The observation can be expanded as
\[
y_k=\langle Y,f_k\rangle=s_k\theta_k+\frac1{\sqrt n}\xi_k,
\]
where \(\xi_k\) are independent standard Gaussian variables. For a cutoff \(K\), define
\[
\hat g_K=\sum_{k\le K}s_k^{-1}y_k e_k.
\]
Then
\[
\E\|\hat g_K-g\|_{L^2}^2
\le \sum_{k>K}\theta_k^2
+\frac1n\sum_{k\le K}s_k^{-2}.
\]
Since \(\|g\|_{H^\alpha}\le1\),
\[
\sum_{k>K}\theta_k^2\le C K^{-2\alpha/d}.
\]
Using \(s_k^{-2}\asymp k^{2/d}\),
\[
\frac1n\sum_{k\le K}s_k^{-2}
\le \frac{C}{n}K^{1+2/d}.
\]
Choose
\[
K\asymp n^{d/(2\alpha+d+2)}.
\]
Then both terms are bounded by
\[
C n^{-\frac{2\alpha}{2\alpha+d+2}}.
\]
Hence the upper bound follows.

\medskip
\noindent\textbf{Lower bound.}
Let \(M\) be a large frequency scale. In a fixed region where \(a(x)\) is bounded away from zero, consider \(N\asymp M^d\) orthonormal wavelet packets \(\psi_{M,l}\), \(l=1,\dots,N\), at frequency scale \(M\). For a wavelet packet normalized in \(L^2\), we have
\[
\|\psi_{M,l}\|_{C^\alpha}\asymp M^{\alpha+d/2},\qquad
\|A\psi_{M,l}\|_{L^2}\asymp M^{-1}.
\]
Let
\[
\epsilon\asymp M^{-\alpha-d/2}.
\]
Then the perturbations \(g_l=\epsilon\psi_{M,l}\) satisfy \(\|g_l\|_{C^\alpha}\lesssim1\), and their \(L^2\)-norm is \(\asymp\epsilon\). Consider the finite family
\[
g_\omega=\sum_{l=1}^N \omega_l g_l,\qquad \omega_l\in\{-1,1\}.
\]
The separation in \(L^2\) between two distinct elements of this family is at least \(\asymp\epsilon\), hence the squared separation is \(\asymp\epsilon^2\). The observed signal \(A g_\omega\) has components of size \(\asymp \epsilon M^{-1}\) in each wavelet direction. The noise level is \(n^{-1/2}\). Thus the signal-to-noise ratio per coefficient is
\[
n^{1/2}\epsilon M^{-1}
\asymp n^{1/2}M^{-\alpha-d/2-1}.
\]
To make the problem statistically hard, choose \(M\) so that this ratio is \(O(1)\), i.e.,
\[
M^{2\alpha+d+2}\asymp n.
\]
Then \(M\asymp n^{1/(2\alpha+d+2)}\). By Fano's inequality (or the standard reduction to a finite Gaussian testing problem; see \cite[Chapter 2]{tsybakov2009}), the minimax \(L^2\)-risk is at least a constant times
\[
N\epsilon^2\asymp M^d\cdot M^{-2\alpha-d}
=M^{-2\alpha}
\asymp n^{-\frac{2\alpha}{2\alpha+d+2}}.
\]
This proves the lower bound.
\end{proof}

\section{Completion of the proof of Theorem~\ref{thm:main}}

By the local linearization in Section~2, the nonlinear model \(Y=u_f+n^{-1/2}\mathbb W\) has, up to constants, the same minimax rate as the linear model \(Y=A g+n^{-1/2}\mathbb W\), with \(g=f-\one\in B_{C^\alpha}(0,r)\). Since \(C^\alpha(D)\) embeds continuously into \(H^\alpha(D)\), the upper bound of Lemma~\ref{lem:linear-rate} applies to the nonlinear problem.

For the lower bound, the wavelet perturbations used in the proof of Lemma~\ref{lem:linear-rate} can be chosen with sufficiently small \(C^\alpha\)-norm and supported away from the zero \(x=0\) of \(a\). For such perturbations, the nonlinear remainder in \eqref{eq:nonlinear-v} is of order \(O(\epsilon^2 M^{-1})\), while the linear signal is of order \(\epsilon M^{-1}\). Since \(\epsilon M\asymp M^{-\alpha-d/2+1}\to0\) as \(M\to\infty\) (because \(\alpha>1\)), the nonlinearity is negligible at the frequency scale \(M\asymp n^{1/(2\alpha+d+2)}\). Hence the same lower bound applies to the nonlinear model.

Combining the two bounds yields
\[
\inf_{\hat f_n}\sup_{f_0\in\mathcal F(\alpha,r)}
\E_{f_0}\|\hat f_n-f_0\|_{L^2(D)}^2
\asymp
n^{-\frac{2\alpha}{2\alpha+d+2}}.
\]
This completes the proof of Theorem~\ref{thm:main}.

\begin{thebibliography}{9}

\bibitem{nickl2019}
R.~Nickl and G.~P.~Paternain.
\newblock On the statistical theory of nonlinear inverse problems.
\newblock \emph{Annals of Statistics}, 47(6):3455--3485, 2019.

\bibitem{tsybakov2009}
A.~B.~Tsybakov.
\newblock \emph{Introduction to Nonparametric Estimation}.
\newblock Springer, 2009.

\bibitem{hormander1985}
L.~Hörmander.
\newblock \emph{The Analysis of Linear Partial Differential Operators III}.
\newblock Springer, 1985.

\end{thebibliography}

\end{document}